\section{Counting sparse roots in one multiplicative disk}\label{sec:local}
Normalize $v_p(p)=1$ on $\overline{\Q}_p$. In the following arguments a finite complete extension contains all points and all scaling constants used. Write $\mathcal O$ for its valuation ring. Define, for a rational $\lambda>1/(p-1)$,
\[
 w_{p,\lambda}(n)=n\lambda-v_p(n!),\qquad w_{p,\lambda}(0)=0.
\]
\begin{lemma}[Local bound]\label{lem:disk}
Suppose $P$ has $k$ nonzero terms, $u\ne0$, and $v_p(b)=\lambda$. The number of distinct roots of $P$ in $u(1+b\mathcal O)$ is at most the largest integer $n$ such that
\[
 w_{p,\lambda}(n)\le\max_{0\le i<k}w_{p,\lambda}(i).
\]
The same bound holds on the corresponding open disk.
\end{lemma}
\begin{proof}
Use the bijection $X=u\exp(bY)$, $Y\in\mathcal O$, whose inverse is $\log(X/u)/b$. For $P(X)=\sum_{j=1}^k c_jX^{e_j}$ with distinct nonnegative integer exponents put $A_m=\sum_j c_j u^{e_j}e_j^m$, using $e_j^0=1$. Then
\[
 P(u\exp(bY))=\sum_{m\ge0} A_m\frac{b^m}{m!}Y^m.
\]
The first $k$ moments are not all zero by the ordinary Vandermonde determinant in the $e_j$. The monic polynomial $\prod_j(T-e_j)$ gives the sequence an integral linear recurrence. The ultrametric inequality therefore gives $|A_m|_p\le M=\max_{i<k}|A_i|_p$ for every $m$. At an index $i<k$ realizing $M$, the absolute coefficient is $M p^{-w_{p,\lambda}(i)}$. Every index past the displayed bound has strictly smaller coefficient. The series converges on the closed disk because $w_{p,\lambda}(m)\to\infty$. Strassmann's theorem bounds the number of zeros by the last index attaining the maximum coefficient norm. This proves the assertion. For rational $\lambda$, a finite extension realizing the valuation of $b$ always exists.
\end{proof}

This argument is related to Lenstra's local fewnomial bounds and to Avenda\~no and Krick's Proposition 3.8. Their proposition supplies the bound $k-1$ in a principal-unit disk if $p>e+k-1$, where $e$ is the ramification index. Retaining the individual factorial valuations in Lemma~\ref{lem:disk} handles small residue characteristic and some large values of $k$.\cite{aven}

\begin{corollary}\label{cor:coarse}
In an unramified local field at an odd $p$, a disk specified by one nonzero residue contains at most $2(k-1)$ distinct roots of a $k$-term polynomial. For fixed $p$ it contains at most $k-1+O_p(\log(2k))$.
\end{corollary}
\begin{proof}
Lift the residue to a unit $u$ and take $b=p$. Write $t=k-1$. Since $v_p(n!)\le n/(p-1)$, we have $w_{p,1}(n)\ge n(p-2)/(p-1)$ and $\max_{i\le t}w_{p,1}(i)\le t$. This proves the first estimate. More precisely, writing $s_p(n)$ for the sum of the base-$p$ digits gives
\[
 w_{p,1}(n)=\frac{(p-2)n+s_p(n)}{p-1}.
\]
For $i\le t$ the digit sum is at most $(p-1)(1+\log_p(\max(1,t)))$. Substituting in the lemma gives the second estimate.
\end{proof}

\paragraph{Counting the conjugates in a certified disk.}
If $H$ is monic and separable, compute $B_i$ by expanding $H(p^a(u+Y))=\sum_i B_iY^i$ for a rational integer $u$ prime to $p$. The smallest index minimizing $v_p(B_i)+i\lambda$ equals the number of roots for which $v_p(p^{-a}X-u)>\lambda$. This is the Newton polygon root-counting rule. To see the endpoint convention, substitute $Y=bZ$, divide by a coefficient of minimum valuation, and reduce. A root strictly inside contributes a factor $Z$, a root on the boundary a factor with nonzero constant term, and an exterior root contributes a constant after normalization. Thus the order at $Z=0$ of the reduction counts precisely the interior roots. Division of the center by its unit part converts the disk to that in Lemma~\ref{lem:disk}. Rational noninteger radii require no computations in a ramified field in the verifier.
